\section{Motivation and significance}
\label{sec:MotivationAndSignificance}

Modern Geographic Information Systems (GIS) and domain-specific data spaces manage vast volumes of multi-attribute spatial, environmental, and observational records~\cite{curry2022data}. Interacting with these platforms typically requires users to navigate complex relational schemas, master rigid query syntax (such as SQL or Lucene), or execute dozens of manual filter selections across multi-level nested web forms~\cite{huang2022natural, dou2024survey}. This introduces substantial cognitive friction, limits exploratory spatial inquiry, and prevents non-technical domain stakeholders from effectively discovering relevant geographic assets.

Recent breakthroughs in Large Language Models (LLMs) and conversational artificial intelligence have created promising opportunities for natural language querying in geospatial domains (GeoAI)~\cite{janowicz2020geoai, mai2023opportunities, li2023geollm}. However, deploying unconstrained LLMs over spatial databases introduces three severe impediments: (i)~\textbf{Database Schema Hallucination}, where models invent non-existent attribute names or relational keys~\cite{dou2024survey, willard2023outlines}; (ii)~\textbf{Non-Deterministic Output Structures}, producing unparseable payloads that break downstream database APIs~\cite{willard2023outlines, pydantic2023}; and (iii)~\textbf{Absence of Geospatial Visual Feedback}, isolating users in text streams without cartographic synchronization~\cite{scheider2021question, huang2022natural}.

To bridge this gap, this paper presents a generic, modular, open-source software architecture connecting a dual-variant Natural Language Understanding (NLU) pipeline, Apache Solr structured spatial indexing, and dynamic Geographic Information System (GIS) map visual synchronization. The architecture is domain-agnostic: by establishing clean functional boundaries between intent parsing, terminology normalization, search index querying, and cartographic rendering, the pipeline can be instantiated across diverse spatial disciplines by modifying only the domain thesaurus and the OpenAPI schema specification.

\begin{table}[h!]
    \centering
    \small
    \setlength{\tabcolsep}{5pt}
    \caption{Comparative analysis of the proposed NLU-Solr-GIS architecture against existing spatial search and conversational AI approaches.}
    \label{tab:tool_comparison}
    \begin{tabularx}{\linewidth}{>{\raggedright\arraybackslash\bfseries}p{3.4cm}>{\raggedright\arraybackslash}X>{\raggedright\arraybackslash}X>{\raggedright\arraybackslash}X}
        \toprule
        \textbf{Dimension} & \textbf{Traditional GIS} & \textbf{Raw LLMs / Text-to-SQL} & \textbf{Proposed Architecture} \\
        \midrule
        Query Modality & Rigid forms / Dropdowns & Free-form natural text & Conversational spatial NLU \\
        Hal\-lu\-ci\-na\-tion Risk & None (Static controls) & High (Invented attributes) & \textbf{Zero (OpenAPI Schema)} \\
        Output Determinism & High & Low (Stochastic output) & \textbf{Guaranteed (Validated JSON)} \\
        GIS Map Feedback & Manual navigation & None (Text-only stream) & \textbf{Real-Time Pulse \& Badges} \\
        Domain Lexicon & Exact string matching & Latent / Unpredictable & \textbf{Pluggable Thesaurus} \\
        Data Sovereignty & Local database & Low (External cloud APIs) & \textbf{100\% On-Premises GPU} \\
        Repro\-duci\-bil\-i\-ty & Live database required & External API key required & \textbf{Standalone Mock Mode} \\
        \bottomrule
    \end{tabularx}
\end{table}

\subsection{Case Study and Validation Testbed: The CRMsDataSpace European Project}
\label{subsec:CaseStudyContext}

To demonstrate the practical applicability of the proposed architecture within a demanding, real-world deployment, the software is instantiated and validated within \textbf{CRMsDataSpace}, a European research project funded by the Research Fund for Coal and Steel (Grant Agreement No.~101216677). 

The objective of the CRMsDataSpace initiative is to establish a Common European Data Space on Critical Raw Materials (CRMs) for the European Green Deal~\cite{european2024crma, blengini2020eu}. The project characterises closed extractive waste facilities, tailings ponds, and spoil heaps across Europe as secondary mineral sources to support EU supply resilience~\cite{vidal2017environmental, lottermoser2010mine}. Within this ecosystem, our software architecture serves as the conversational front-end of the Minimum Viable Data Space (MVDS), allowing mining engineers, geologists, and policymakers to query multi-criteria spatial and mineralogical records in natural language with synchronized cartographic feedback. The CRMs domain acts as an exemplary case study highlighting the versatility and robustness of the generic software pattern.
